\chapter[An Overnight Scientific Revolution?]{Did a Few Geniuses Really Make the Scientific Revolution Overnight?}
\section{A Glass Prism}
Pick up something small enough to hold in one hand: a triangular piece of glass with polished, shining edges. It is a prism.

You have probably seen one. In primary-school science, a teacher held it at a window and passed sunlight through it, spreading a little rainbow across the wall: red, orange, yellow, green, blue, indigo and violet, as though a palette had spilled into the light. How can transparent glass divide white light into seven colours? Is white light ``pure'', or seven colours mixed together in disguise?

Over three centuries ago, a young man in his twenties pondered the same question. Around 1665, plague closed Cambridge University, sending the student Isaac Newton home to the countryside. He bored a small hole in his bedroom shutter and passed the entering sunlight through a prism. The prevailing explanation held that prisms coloured light: glass contaminated pure white light, producing colours. Newton added a step others had not considered. He isolated a red beam from the rainbow and passed it through another prism. If prisms supplied colour, red should deepen. Instead, red remained red, unchanged. The conclusion reversed: prisms did not produce colours, but separated colours already mixed within white light. White was the mixture; colours were the originals.

The experiment entered every physics textbook, usually occupying half a page with ``Newton discovered the dispersion of light''. Weigh the prism in your palm, however, and questions emerge that textbooks leave unanswered. Where did the glass come from? Why did he think of adding a second prism? Arab scholars and Oxford friars had performed related optical experiments centuries earlier; why did this one count? Above all, we habitually tell this as a flash of genius. Could a piece of glass, a hole and an idle holiday really lever a ``Scientific Revolution'' into existence?

That is our subject. Starting with the prism, we shall ask whether the Scientific Revolution was a miraculous night for a handful of geniuses or a slow fire stirred by countless hands over centuries.

\section{An Observatory Castle on Hven}
Come into a scene.

The time is a winter night in 1585, far below freezing, with Baltic winds cutting like knives. The place is Hven, a small island off Denmark. The Danish king owns it and has granted the entire island, tenants included, to the nobleman Tycho Brahe. With royal money, Tycho builds one of Europe's most extravagant structures: neither palace nor church, but an astronomical castle he calls Uraniborg, the ``castle of the heavens''.

Follow a young assistant up its tower. Your nose first reveals an unusual household: pulp and ink scent the stairs. Dissatisfied with outside paper and slow printing, Tycho established his own paper mill and press, letting observations become printed pages on the spot. Wind pours across the tower top. Tonight is cloudless, improbably crowded with stars. Tycho stands beside a huge brass quadrant: essentially a great fan-shaped plate precisely graduated to measure stellar altitude. The metal clings coldly to hands, its divisions fine as hairs. Torchlight gives the tip of his nose an unnatural metallic glint. A youthful duel cut away part of his nose; he wears a metal replacement for the rest of his life.

The assistant's task is astonishingly monotonous: Tycho calls numbers; you record them. ``Mars, altitude \ldots'' Night after night, year after year, for twenty years. No telescope: its invention is still over two decades away. Only naked eyes, brass instruments and near-obsessive care. Tycho reduces instrumental error beyond his predecessors' dreams. At his death, the accumulated records form humanity's cleanest, most substantial archive of the sky.

Notice things more important than any genius: royal money, without which there is no Uraniborg; papermaking and printing, without which unshared data might as well not exist; brassworkers' skill, without accurate graduations even Tycho's genius is wasted; and the anonymous assistant recording numbers through the night, whose name history did not preserve but whose hands carried every figure.

What of Tycho himself? His twenty years of observation defended a universe we now consider wrong. Rejecting Earth's orbit around the Sun, he proposed a compromise: the Sun circled Earth, while the other planets circled the Sun. The world's most diligent, precise observer believed a model we would reject tonight. That fact contains our question. If genius's eyes could make a Scientific Revolution, why did the finest pair look for twenty years and ``see wrongly''?

\section{One Night or a Century and a Half?}
Set out the conflict without rushing to answer.

You probably grew up with a heroic account. Ignorant medieval people believed Earth stood at the universe's centre beneath perfect, sacred heavens. Geniuses then appeared one by one: Copernicus boldly proposed heliocentrism; Bruno died for it; Galileo raised a telescope and overturned the old world; an apple struck Newton beneath a tree, universal gravitation was born, and the revolution was complete. Like dominoes, each genius knocked down another part of the old world. Easy to tell, remember and reproduce in examinations.

Spread the same events across their dates and the ``night'' looks different:

In 1543, the Polish cleric Copernicus published On the Revolutions of the Heavenly Spheres, proposing Earth's orbit around the Sun. He was gravely ill at publication, and few believed him for decades. In 1572, a new star appeared. Tycho measured its lack of parallax, a term we shall explain, showing change in supposedly unchanging heavens. In 1600, Gilbert published On the Magnet, treating the compass needle's direction as a serious natural question for the first time. In 1609, Kepler's New Astronomy announced elliptical planetary orbits, while Galileo pointed a telescope skywards. Bacon's New Organon in 1620 advocated systematic observation and experiment. Descartes published the Discourse on Method in 1637. In 1660, London enthusiasts established the Royal Society; within years came one of the earliest scientific journals, Philosophical Transactions. Newton's Mathematical Principles of Natural Philosophy appeared in 1687, unifying heavenly planets and earthly apples under mathematical laws.

From 1543 to 1687: 144 years, a four-generation relay. No individual leg looked much like a revolution. Copernicus made few new observations, relying largely on ancient data. Tycho observed for twenty years to defend a hybrid geocentric model. Kepler used Tycho's posthumous data. Galileo experimented, but his sharpest weapons were mathematics and imagined ideal conditions. Newton openly built on everyone's results. No leg lasted one night; none can even be detached as a revolution on its own. Remove a link and the next runner cannot receive the baton.

The real question emerges: \textbf{how should we understand this transformation of humanity's destiny?} Was it the triumph of exceptional minds, without Newton leaving humanity in darkness for centuries? Or an inevitable current generated jointly by economics, instrument-making, printing and oceanic trade, in which someone else would have formulated gravitation around the same time? The Chinese wordplay substitutes Ma-dun or Yang-dun, ``horse-Newton'' or ``sheep-Newton'', for Niu-dun, Newton's Chinese name, whose first syllable also means ``ox''. Or are both questions themselves mistaken?

Before reading on, cast an inward vote: without Newton, would someone still have discovered universal gravitation around 1700?

\section[Opponents and Other Observers]{Dissenting Votes inside, Stargazers Outside}
Heroic stories casually turn opponents into fools and villains. Undo that move by fully presenting the reasons of people who rejected heliocentrism. Only after hearing their strongest case can you weigh the revolution's achievement.

\textbf{A word for the old school.} Around 1600, an educated European could reject heliocentrism with at least four strong cards.

First, common sense and sensory experience. If Earth raced at extraordinary speed, why did a stone dropped from a tower land beside its base, not far to the west? Why were birds flying vertically not left behind? Unless heliocentrism explained our inability to feel Earth's motion, it contradicted daily bodily experience.

Second, accuracy. You may imagine Copernicus's system immediately calculated better. Quite the opposite. Holding to the ancient requirement of perfect circular celestial motion, he needed numerous small circles to fit observations. His planetary positions were no more accurate than Ptolemy's. Why should a theory win if it predicts no better while demanding belief in a racing Earth?

Third, strongest of all, stellar parallax. If Earth travelled a large orbit, viewing a star six months apart from opposite ends should shift its apparent position slightly, like a finger jumping against the background when you alternately close each eye. Observers searched for decades and found no such shift. Heliocentrists could only reply that stars were too distant for measurable parallax. Correct today, this then amounted to asking belief in an unsupported assumption. Direct observation of Earth's orbit came only in 1838, when German astronomer Bessel measured it, 295 years after Copernicus's book.

Fourth, interpretative tradition. Biblical interpretation then read certain passages literally. Recently shaken by the Reformation, the Church was especially sensitive to freely reinterpreting Scripture. Galileo was tried by the Inquisition in 1633 and sentenced to lifelong house arrest: that is fact. Reducing it to ``an ignorant Church persecutes science'', however, omits the background. He had enjoyed good relations with senior churchmen and remained devout. The trial entangled a political disaster: placing arguments of a pope who had been his friend into a foolish character's mouth. History is a tangle, not a cartoon.

The opposing votes in that room were therefore not stupid votes. This matters enormously. The revolution did not mean clever people defeating fools. Over decades, the new system accumulated things the old could not offer: better predictions, easier calculations and new telescopic facts. Its hard-won victory shows the achievement was real.

\textbf{Now widen the view beyond the room.} While Europeans argued over geocentrism and heliocentrism, people elsewhere quietly observed the stars, often very well.

In 1259, an Ilkhanate-funded observatory arose at Maragha in present-day Iran, directed by Nasir al-Din al-Tusi. He devised an ingenious mathematical arrangement, later called the Tusi couple, combining circular motions to approximate straight-line oscillation. A strikingly similar device appeared in Copernicus's book three centuries later. Scholars still debate the precise transmission, but broadly agree that European astronomy's mathematical toolbox contained Islamic components.

In China, Yuan astronomer Guo Shoujing directed extensive measurements and introduced the Season-Granting Calendar in 1280. Its tropical-year value matched the Gregorian calendar's three centuries later. At Beijing's Imperial Astronomical Bureau, observation, recording and calendar-making belonged to a state undertaking almost two millennia old. Notice the same combination of peers, state patronage and precise instruments as Uraniborg. Astronomy was never Europe's solo performance. Traditions answered somewhat different questions: Chinese calendars served imperial legitimacy and agriculture; Islamic observatories served calendars, astrology and religious times; later European astronomy increasingly pursued a theoretical question, the true geometry of planetary motion.

Quieter voices rarely enter heroic histories. In the seventeenth-century Andes, European missionaries and physicians learned to treat malaria with cinchona bark, knowledge Indigenous people had held for generations. Jesuits brought it to Europe, where it became a major medicine for two centuries. In 1769, Captain Cook's expedition sailed to Tahiti chiefly to observe Venus's transit, joining worldwide measurements to calculate the Earth--Sun distance. Naturalist Joseph Banks collected boxes of plants. Global shipping connected European scholars to observation sites and knowledge stores worldwide. The same routes carried colonial conquest, plunder and slave trading; that account cannot be omitted. Science's global data network was empire's expansion network. Knowledge acquired new names in European print; its original suppliers usually remained unnamed.

Inside the room are at least two positions, with strong reasons on the old system's side. Outside are at least two neglected groups: other civilisations' astronomical traditions and anonymous knowledge providers within colonial networks. A story about a few European geniuses is growing too small to hold them all.

\section[Thinking Bridge: Lengthen the Timeline]{Thinking Bridge: Put ``Revolution'' on a Long Timeline}
For disputes over overnight change, historians have a clumsy but effective tool: lengthen the timeline. Move a story's beginning backwards until it is not yet recognisably itself, then count the people and events between. Its success at dismantling overnight myths may surprise you. It works for the Renaissance, Industrial Revolution, internet boom and any tale beginning ``One day, someone suddenly \ldots''.

Practise with a question: in what year was universal gravitation discovered?

The heroic answer: 1666, the year the apple fell.

First, trace backwards. Newton's Principia appeared in 1687, twenty-one years after the apple. Between them he wrote optical papers for the Royal Society, quarrelled with Hooke so bitterly he nearly renounced publication, experimented in alchemy and studied biblical chronology. Second, trace further. In 1684, Halley visited Cambridge to ask the shape of a planetary orbit under an inverse-square attraction. Newton immediately answered: an ellipse. Notice that the question already presupposed an inverse-square idea in circulation. Hooke and Halley had considered it; what was missing was a mathematical derivation of the orbit. Third, trace mathematical components: Kepler's ellipse in 1609, Tycho's Mars data from 1576--1601, brassmaking precision and twenty years of Danish royal funding. Fourth, trace conceptual components: Galileo's falling-body studies partly prepared the idea of common earthly and heavenly laws; nature's mathematical language reaches further back to scholastics and Arabic algebra.

By now, ``1666'' is strained. More honestly, the discovery moment stood at the end of a conveyor belt over a century long. Newton performed the final operation, the most crucial, but still the final one.

While here, unpack another school formula: the five-step scientific method---question, hypothesis, experiment, data analysis, conclusion, repeated. Memorising it is harmless. Comparing it with this timeline reveals an embarrassment: \textbf{almost none of the key figures followed those five steps}. Copernicus mostly mathematically beautified old data. Tycho observed while rejecting heliocentrism to his death, appearing under the formula as an incomplete scientist collecting data without hypotheses. Kepler calculated from others' data for a decade, trying dozens of geometries like an accountant refusing defeat. Galileo relied heavily on idealisation---no air resistance, perfectly smooth planes---impossible in real experiments but possible in thought. Newton simply declared that he framed no hypotheses, first completing the mathematics.

They were not doing science wrongly. The scientific method is not an inherited secret recipe but a toolbox: observation, mathematisation, experiment, instruments and peer review, used in different orders. Perhaps the common minimum is simply this: \textbf{eventually you must produce something others can check}---data, a derivation or a reproducible phenomenon. Shorter than the five steps, this requirement is harder: it transfers judgement from genius's confidence to peers' eyes. We shall see how important that transfer from personal authority to communal scrutiny becomes.

\section{On Whose Shoulders?}
Read a short primary source.

Early in 1676, Newton wrote to Hooke. Their relationship was delicate after public optical disputes. The letter contained words that travelled worldwide:

\begin{quote}
If I have seen further, it is by standing on the shoulders of giants. (Newton's 1676 letter to Robert Hooke.)
\end{quote}
Nearly every secondary classroom has displayed this as Newtonian modesty. Knowing the context reveals layers. Hooke was among Newton's fiercest priority rivals, so it also contains a polite distancing: ``My achievements did not come from you.'' Better still, the saying has an older source. Twelfth-century French scholar Bernard of Chartres described modern people as dwarfs on giants' shoulders, seeing farther through their elevation. John of Salisbury recorded it, and it circulated for centuries. Even saying ``on the shoulders of giants'' stands on predecessors. A maxim celebrating individual genius is itself evidence of collective relay: an excellent irony.

Now a paraphrase. In 1610, Galileo published the little Starry Messenger, reporting his homemade telescope's observations. The Moon was not smooth and polished, but rough, with elevations and depressions resembling Earth's mountains and valleys. Four little stars beside Jupiter changed position nightly while circling it. The Milky Way's pale band resolved into countless stars.

Read these sources together. Starry Messenger challenged a belief not seriously shaken for two millennia: heaven and Earth were different kinds of world, one perfect, eternal and smooth, the other rough and perishable. Lunar mountains made heaven a larger Earth; Jupiter's moons revealed centres of motion other than Earth. Yet every revolutionary sentence depended on three impersonal things: Dutch spectacle-makers' combinations of lenses, which Galileo improved after hearing of them; Venetian glassworkers' grinding skills; and printing. Within months, the book spread across Europe and colleagues made telescopes to check. Some accepted what they saw. Others refused even to look, dismissing lens-generated deception. The test was whether others could see it too. Uraniborg's press, Royal Society journals and quarrelsome correspondence networks were the revolution's infrastructure. Geniuses supplied sentences; communities decided which survived.

\section{Three Accounts of One Revolution}
We can now offer genuinely different explanations of the same facts. Separate four commonly confused categories. What happened---European astronomy and physics changed foundations between 1543 and 1687---is largely undisputed evidence. Why it happened is where explanations contend. Participants' understanding, often not of a revolution at all, with Tycho repairing an old system, is their self-description. Whether the transformation was fortunate or also seeded arrogance is our evaluation. Here we compare the second category.

\textbf{Explanation one: ideas move themselves---internalism.}

Here the revolution solved successive intellectual problems. Copernicus sought elegance in circular systems, Kepler sought compatibility between Copernicus and Tycho's data, and Newton sought to turn Kepler's empirical laws into mathematical necessities. Each step answered cracks left by the last; society, economics and navigation were background noise. This respects intellectual work's texture. The eight-minute discrepancy in Tycho's data genuinely tormented Kepler for years. No social factor can replace an individual's obsession with a problem and refusal to let theory escape data. Yet it cannot explain why the sixteenth and seventeenth centuries. Ptolemy's cracks had existed for a millennium; Arab and Byzantine scholars knew them. Why did they become great fissures then, on those shores?

\textbf{Explanation two: social needs propel science---externalism.}

From the sixteenth century, European ocean travel demanded accurate star tables and longitude; artillery required ballistics; mine drainage required pumps and vacuum research; commerce and banking needed arithmetic and accounts; emerging states wanted maps, calendars and statistics. Even Copernicus's book ripened through the Church's practical calendar-reform needs. Mid-fifteenth-century printing first enabled cheap reproduction, comparison and correction. In a manuscript age, a datum copied wrongly three times could never recover. Newton thus becomes the inevitable product of ready demand, technology and communication. Without him, Hooke, Halley and Leibniz's cohort would assemble broadly similar results. Newton and Leibniz's near-simultaneous independent calculus provides strong evidence of an idea whose time had come. This sharp account explains timing, but needs do not generate answers automatically. Navigation wanted longitude for two centuries; generations of instrument-makers and astronomers supplied painstaking work, not a slogan about social demand. Externalism also risks reducing science to society's servant, missing apparently useless obsessions such as Kepler's decade devoted to Mars.

\textbf{Explanation three: craftspeople, instruments and new knowledge communities---practice.}

Move from studies to workshops. The central figures were surrounded by artisans: Galileo observed practical knowledge of force and materials in Venice's Arsenal; spectacle-makers produced telescopes; instrument-makers built vacuum pumps; clockmakers enabled precision timing. The novelty was not one theory, but a combination: \textbf{scholars' mathematics and artisans' skills first joined hands on a large scale}, while societies, journals and correspondence tested and circulated results. Older universities ranked manual knowledge low and intellectual knowledge high. Revolution occurred where that wall fell. This fills gaps between the first accounts, explaining both timing---maturing lenses, clocks and printing---and texture: craftsmanship and data challenging theory. Its weakness is flattening differences. Newton and Hooke become difficult to distinguish, yet Newton alone accomplished the mathematical synthesis others did not. Erasing individual contribution altogether is another distortion.

Each account holds part of the truth. This is no evasive compromise: their methodological differences are real. Do we seek causes in minds, social structures or the meeting point of people and things?

\section[Further Challenge: Paradigms]{Further Challenge: Paradigms and Missing Disciplinary Boundaries}
Now introduce the chapter's weightiest theory. In 1962, Thomas Kuhn published the slim Structure of Scientific Revolutions, probably among the twentieth century's most cited scholarly books. He asked our question: how does science change?

Start with three terms. First, \textbf{paradigm}. Mature science is not isolated individuals: a community shares defaults about worthwhile questions, good answers, instruments and paper-writing. Ptolemaic astronomy and Newtonian mechanics are paradigms. Like spectacles, they clarify vision while limiting it to their lenses. Second, \textbf{normal science}. Once a paradigm forms, most work is puzzle-solving within it, filling empty squares rather than overturning the board. Tycho's twenty years become exceptionally good puzzle-solving in an old paradigm, indispensable for exposing its problems. Third, \textbf{anomalies and crisis}. Some nails refuse to enter the board: an eight-minute discrepancy, unmeasurable stellar parallax, Mercury's slight orbital precession. Most anomalies receive patches. Occasionally one grows until the community doubts its spectacles. Then comes revolution: a new paradigm.

The theory has two faces here. Its sharp face explains how the best eyes could ``see wrongly'' for twenty years: Tycho perfected work within an old paradigm. It explains Copernicus's initially sparse following: young paradigms are often rougher than established ones. They win not through immediate superiority but by gradually accumulating what the old cannot provide. It also punctures arrogance. Aristotle's physics was not two millennia of mindless nonsense, but a once-powerful paradigm supporting generations of normal science. Calling it mere ignorance invites future successors to dismiss us likewise.

Weigh Kuhn's limits just as firmly. ``Paradigm'' is overused, especially in business talk, where every grand claim becomes a paradigm shift. A deeper objection asks: if paradigms are separated by incompatible language, as Kuhn sometimes suggests, does science still accumulate and progress? Kepler used Tycho's data; Newton used them; today's orbital calculations still do. Data seem to outlive theories. This bright lamp misses corners, especially why later paradigms genuinely predict better.

History supplies another point, not written into Kuhn's book, answering a further question: \textbf{the word ``science'' did not yet exist in that age}. Copernicus, Kepler, Galileo and Newton called their work natural philosophy, reasoning philosophically about nature. The English ``scientist'' was coined only in 1833. Modern disciplinary walls had not been built. Kepler calculated elliptical orbits while earning his living from horoscopes, sincerely believing celestial arrangements related to human affairs. Newton devoted far more ink and paper to alchemy and biblical chronology than to the Principia. After examining his alchemical manuscripts, Keynes famously called him, in effect, not the first modern man of reason, but the last magician. This was no sneer. Alchemy was serious investigation of material change, not fraudulent science's opposite; Boyle, another founder of chemistry, shared it. Astrology and astronomy were related crafts, both practised at Uraniborg. We retrospectively paste the Principia into a science album and file alchemy under eccentricity, partitioning old rooms with modern walls. The revolution was not simply correct knowledge defeating error. It was a long \textbf{division of a household}: establishing which questions counted. That boundary-making was among the era's hardest and most easily forgotten struggles.

Correct one scene familiar to almost every Chinese child: Galileo dropping two iron balls from Pisa's leaning tower. Its sole source is his pupil Viviani's biography, written years after Galileo's death. Galileo's extensive manuscripts and letters never record it, nor does a contemporary Pisan eyewitness. Most scholars regard it as an attractive teaching legend. He did study falling bodies through thought experiments and inclined planes, but the public drop supposedly demolishing an old worldview before thousands probably never happened. The story itself exhibits our argument: later generations hunger to compress long change into one genius's dramatic instant.

\section{Today's Genius Factory}
This four-hundred-year-old question repeats daily in your news feed.

Count heroic narratives in science headlines: a team's major breakthrough, a teenage genius solving a century-old problem, a paper changing the world overnight. Journalism faces a real difficulty. Collaboration is prolonged, anonymous and visually undramatic; editors need headlines. Science communication keeps mass-producing apple moments: hundreds of people working for a decade become one name and a grand statement. The result is a pervasive illusion that progress means waiting for geniuses to be born. Parents demand the next Einstein; young people doubt themselves for lacking brilliant ideas. Nobody tells them that modern science most needs the ability to do something tedious correctly alongside a roomful of ordinary people.

Actual science differs from the heroic version even more than four centuries ago. A large particle-physics paper may list hundreds or thousands of authors from institutes in dozens of countries. Yet Nobel rules still limit each annual prize to three people, institutionally endorsing the few-geniuses story and annually disappointing a fourth person and whole teams. Scientists debate this endlessly: rewards inhabit Newton's portrait while research resembles Uraniborg multiplied ten thousandfold.

Communal scrutiny continues too. Preprint sites publish before review, letting worldwide peers identify errors overnight. During major epidemics, data reach laboratories worldwide within weeks: the Royal Society's correspondence network at light speed. Communities have ailments, however. The recent reproducibility crisis includes celebrated psychological and medical experiments failing when others repeat them. That returns us to our minimum: something others can check. When checking fails, more papers are only more print. Fringe-science claims reveal another side: outsiders announcing the overthrow of relativity often lack not confidence but patience for communal scrutiny. Gatekeeping can wrongly exclude unconventional minds; without standards, however, science becomes a contest in volume.

Then comes artificial intelligence. In protein-structure prediction, drug-candidate screening and astronomical data analysis, machines are already central performers in some stages rather than assistants. Tycho would instantly recognise, but could not answer, the question: if machines observe and models calculate, whose eyes produced the record? Spectacle-makers once changed who could observe; algorithms now rewrite the question. The Scientific Revolution has no final script, only repeated rearrangements of instruments, data, mathematics, communities and the people between them, sometimes great, sometimes small.

\section[Your Turn: Without the Heroes]{Your Turn: What Remains after Burning the Roll of Heroes?}
Return to the prism.

Imagine that tomorrow all histories of science become anonymous editions. Remove Newton, Einstein and every personal name. Write only: a community in a particular year, using certain instruments, data and disputes, advanced understanding of gravity. Attach each discovery's long prehistory and full author list. Would this make the world more honest?

Before concluding, weigh both sides again.

The case for change holds much of this chapter. Heroic accounts systematically mislead, erasing assistants, artisans, patrons, other civilisations and respectable old-paradigm opponents. They make genius seem science's entry ticket instead of rigour, patience and collaboration. They replace its true judge, public testing, with personal glamour. An inaccurate map remains harmful however beautiful.

Now make the opposing case fully. First, stories open hearts: countless scientists recall a childhood story drawing them towards a laboratory. A map full of unnamed communities may be so accurate nobody feels brave enough to set out. Second, names anchor responsibility. ``Newton's system'' identifies someone answerable for each assertion; ``community'' can be too smooth to hold anyone accountable. Third, hardest perhaps, removing individuality is itself distortion. Newton was the final operation on a conveyor belt, but replacing that operation really changes the product. Leibniz invented calculus without writing the Principia. Timelines dilute genius; they do not evaporate it.

Should an honest history preserve heroes or dissolve them? Can genius stories be the best advertisements and worst maps simultaneously? If so, how will you hold both truths? If not, which will you sacrifice?

There is no standard answer. Notice, however, that answering already uses all our tools: timelines, evidence chains, paradigms and suspicion of any single explanation. From Copernicus's book to the one in your hands, the conveyor belt has never stopped. Who performs its next operation depends partly on what you think after closing this book tonight.
